\section{Slutsatser och avslutande diskussion}

\subsection*{Vad jag framför allt har lärt mig}

Den viktigaste lärdomen från LIA 2 är att ett system inte är färdigt bara för att det går att
deploya. Det behöver också vara säkert, stabilt, lätt att övervaka och möjligt att förbättra utan
att orsaka onödiga avbrott.

Det är en viktig skillnad mot många skolprojekt. Under utbildningen bygger man ofta nya
lösningar. Under LIA 2 lärde jag mig vad det innebär att ta ansvar för ett system som redan
används och som andra är beroende av.

\subsection*{Om jag nådde mina mål}

Jag tycker att jag nådde mina mål i hög grad. Jag arbetade med Kubernetes, Terraform,
Workload Identity, automatiserad deployment och förbättring av CI/CD. Jag lyckades också
förstå hur dessa delar hänger ihop i praktiken.

Det som överraskade mig mest var att jag inte bara lärde mig nya verktyg, utan också hur viktigt
det är att tänka på helheten. Driftstabilitet påverkar förtroendet för systemet. Säkerhet måste
byggas in från början. CI/CD påverkar hur snabbt och tryggt ett team kan arbeta.

\subsection*{Det jag tar med mig vidare}

Det finns tre saker som jag särskilt tar med mig från LIA 2.

För det första har jag lärt mig att säkerhet inte är något man lägger till i slutet. Den måste vara
en del av lösningen från början.

För det andra har jag förstått att tydlig kommunikation är en viktig del av det tekniska arbetet.
Bra dokumentation, tydliga pull requests och raka statusuppdateringar gör att teamet arbetar
bättre tillsammans.

För det tredje har jag utvecklat ett större ansvarstagande. Jag ser nu tydligare att det inte räcker
att bara lösa det som står i uppgiften. Man behöver också se hur den egna lösningen påverkar
helheten.

\subsection*{Avslutning}

Två LIA-perioder på Insighta har gett mig en stark praktisk grund. Jag känner mig nu tryggare i
rollen som blivande DevOps Engineer och mer säker på vilken typ av arbete jag vill fortsätta med
efter utbildningen.